%--------------------------------------------------------------------
\subsection{Periodic RT-TDDFT and the Velocity Gauge}
\label{sec:periodic-velocity-gauge}
%--------------------------------------------------------------------

The starting point for the present periodic RT-TDDFT formulation is
the time-dependent Kohn--Sham equation in the presence of an external
electromagnetic perturbation,
\begin{equation}
  \imath\hbar \frac{\partial}{\partial t}\Psi_j(\mathbf r,t)
  =
  \hat H(t)\Psi_j(\mathbf r,t),
  \label{eq:tdks_em_general}
\end{equation}
where
\begin{equation}
  \hat H(t)
  =
  \frac{1}{2m_e}
  \left[
    \hat{\mathbf p}
    -
    \frac{q_e}{c}\mathbf A(\mathbf r,t)
  \right]^2
  + q_e\Phi(\mathbf r,t)
  + \hat V_{\rm KS}[\rho](t).
  \label{eq:tdks_minimal_coupling}
\end{equation}
Here $q_e=-e$ is the electron charge, with $e>0$, $m_e$ is the
electron mass, $c$ is the speed of light, and
\begin{equation}
  \hat{\mathbf p}=-\imath\hbar\nabla
  \label{eq:canonical_momentum}
\end{equation}
is the canonical momentum operator.  Thus, for an electron, the vector
potential enters the kinetic momentum as $\hat{\mathbf p}+e\mathbf A/c$.
In the norm-conserving pseudopotential formulation used here, the
Kohn--Sham potential is decomposed as
\begin{equation}
  \hat V_{\rm KS}[\rho](t)
  =
  \hat V_{\rm ion}^{\rm loc}
  + \hat V_{\rm NL}
  + \hat V_{\rm H}[\rho](t)
  + \hat V_{\rm xc}[\rho](t),
  \label{eq:vks_decomposition}
\end{equation}
where $\hat V_{\rm ion}^{\rm loc}$ is the local ionic pseudopotential,
$\hat V_{\rm NL}$ is the nonlocal pseudopotential operator,
$\hat V_{\rm H}$ is the Hartree potential, and $\hat V_{\rm xc}$ is the
exchange-correlation potential.

The physical electric and magnetic fields are obtained from the
electromagnetic potentials according to
\begin{equation}
  \mathbf E(\mathbf r,t)
  =
  -\nabla\Phi(\mathbf r,t)
  -
  \frac{1}{c}
  \frac{\partial\mathbf A(\mathbf r,t)}{\partial t},
  \qquad
  \mathbf B(\mathbf r,t)
  =
  \nabla\times\mathbf A(\mathbf r,t).
  \label{eq:fields_from_potentials}
\end{equation}
Different choices of $\Phi$ and $\mathbf A$ related by a gauge
transformation describe the same electromagnetic fields.  This gauge
freedom should be distinguished from the long-wavelength approximation
used for optical response.  In the present work we consider the
homogeneous-field limit, in which the spatial variation of the external
field over the simulation cell is neglected.  This approximation
corresponds to the optical, $\mathbf q\rightarrow 0$, limit and leads to
vertical transitions at fixed crystal momentum.  The
$\mathbf q\rightarrow 0$ restriction therefore arises from the
homogeneous-field approximation, not from the choice of velocity gauge.

For finite systems, a spatially uniform electric field is commonly
represented in the length gauge by setting
\begin{equation}
  \mathbf A_L(\mathbf r,t)=0,
  \qquad
  \Phi_L(\mathbf r,t)=-\mathbf r\cdot\mathbf E(t).
  \label{eq:length_gauge_potentials}
\end{equation}
The corresponding scalar-potential energy for an electron is
$q_e\Phi_L(\mathbf r,t)=e\,\mathbf r\cdot\mathbf E(t)$.  This
representation is well suited to finite molecular systems, where the
time-dependent dipole moment is a well-defined observable.  Under
periodic boundary conditions, however, the position operator is not
compatible with Born--von Karman periodicity: $\mathbf r$ and
$\mathbf r+\mathbf L$ represent equivalent points of the crystal but
give different values of the scalar potential in
Eq.~\eqref{eq:length_gauge_potentials}.  The length-gauge scalar
potential is therefore not a periodic operator.

The same homogeneous electric field can be represented in the velocity
gauge by a gauge transformation.  For a scalar gauge function
$\chi(\mathbf r,t)$,
\begin{equation}
  \mathbf A'(\mathbf r,t)
  =
  \mathbf A(\mathbf r,t)+\nabla\chi(\mathbf r,t),
  \qquad
  \Phi'(\mathbf r,t)
  =
  \Phi(\mathbf r,t)
  -
  \frac{1}{c}
  \frac{\partial\chi(\mathbf r,t)}{\partial t},
  \label{eq:gauge_transformation_potentials}
\end{equation}
and the electronic wavefunction transforms as
\begin{equation}
  \Psi'(\mathbf r,t)
  =
  \exp\!\left[
    \frac{\imath q_e}{\hbar c}\chi(\mathbf r,t)
  \right]
  \Psi(\mathbf r,t).
  \label{eq:gauge_transformation_wavefunction}
\end{equation}
Starting from the length-gauge potentials in
Eq.~\eqref{eq:length_gauge_potentials}, choose
\begin{equation}
  \chi(\mathbf r,t)=\mathbf r\cdot\mathbf A(t),
  \qquad
  \frac{d\mathbf A(t)}{dt}=-c\mathbf E(t).
  \label{eq:goppert_mayer_chi}
\end{equation}
This gives
\begin{equation}
  \Phi_V(\mathbf r,t)=0,
  \qquad
  \mathbf A_V(\mathbf r,t)=\mathbf A(t),
  \qquad
  \mathbf E(t)=-\frac{1}{c}\frac{d\mathbf A(t)}{dt}.
  \label{eq:velocity_gauge_potentials}
\end{equation}
For an electron, the corresponding unitary transformation is
\begin{equation}
  \hat U(t)
  =
  \exp\!\left[
    -\frac{\imath e}{\hbar c}
    \mathbf r\cdot\mathbf A(t)
  \right],
  \qquad
  \Psi_V(\mathbf r,t)=\hat U(t)\Psi_L(\mathbf r,t).
  \label{eq:goppert_mayer_unitary}
\end{equation}
Because $\mathbf A(t)$ is spatially uniform,
$\nabla\times\mathbf A(t)=0$; the homogeneous velocity-gauge field is
therefore the gauge-equivalent representation of the electric-dipole
field, not a full finite-$\mathbf q$ transverse electromagnetic wave.
Finite-$\mathbf q$ response functions lie outside the homogeneous-field
formulation considered here.

This relation also fixes the formal connection between a conventional
broadband electric-field kick and the corresponding velocity-gauge vector
potential.  For an impulsive perturbation,
\begin{equation}
  \mathbf E(t)=\mathbf E_0\delta(t),
  \label{eq:delta_electric_kick}
\end{equation}
integration of Eq.~\eqref{eq:velocity_gauge_potentials} gives, for
$\mathbf A(t<0)=0$ under the sign conventions established above,
\begin{equation}
  \mathbf A(t)=-c\mathbf E_0\Theta(t).
  \label{eq:step_vector_potential}
\end{equation}
The particular impulsive realization used in RMG is described separately
in Sec.~\ref{sec:implementation}.

In the formal velocity-gauge representation, the periodic TDKS
Hamiltonian for an electron is
\begin{equation}
  \hat H_V(t)
  =
  \frac{1}{2m_e}
  \left[
    \hat{\mathbf p}
    -
    \frac{q_e}{c}\mathbf A(t)
  \right]^2
  + \hat V_{\rm loc}[\rho](t)
  + \hat V_{\rm NL}^{V}(t).
  \label{eq:velocity_gauge_hamiltonian}
\end{equation}
Here
\begin{equation}
  \hat V_{\rm loc}[\rho](t)
  =
  \hat V_{\rm ion}^{\rm loc}
  + \hat V_{\rm H}[\rho](t)
  + \hat V_{\rm xc}[\rho](t)
  \label{eq:vloc_definition}
\end{equation}
contains all local terms, while the nonlocal pseudopotential is written
separately.  Since $q_e=-e$, the kinetic momentum in
Eq.~\eqref{eq:velocity_gauge_hamiltonian} is
$\hat{\mathbf p}+e\mathbf A(t)/c$ for an electron.  If a nonlocal
pseudopotential is present, it must be transformed consistently with the
same unitary gauge transformation,
\begin{equation}
  \hat V_{\rm NL}^{V}(t)
  =
  \hat U(t)\hat V_{\rm NL}\hat U^\dagger(t),
  \qquad
  \hat U(t)
  =
  \exp\!\left[
    \frac{\imath q_e}{\hbar c}
    \mathbf r\cdot\mathbf A(t)
  \right].
  \label{eq:gauge_transformed_nlpp}
\end{equation}
The detailed consequences of Eq.~\eqref{eq:gauge_transformed_nlpp} for
the current operator are discussed in Sec.~\ref{sec:current-nlpp}.

For a periodic system, the time-dependent Kohn--Sham orbitals are
written in Bloch form,
\begin{equation}
  \Psi_{n\mathbf k}(\mathbf r,t)
  =
  e^{\imath\mathbf k\cdot\mathbf r}
  u_{n\mathbf k}(\mathbf r,t),
  \label{eq:bloch_orbital_td}
\end{equation}
where $u_{n\mathbf k}(\mathbf r,t)$ has the periodicity of the unit
cell.  Acting on a Bloch orbital, the canonical momentum becomes
\begin{equation}
  \hat{\mathbf p}\Psi_{n\mathbf k}
  =
  e^{\imath\mathbf k\cdot\mathbf r}
  \left(
    -\imath\hbar\nabla+\hbar\mathbf k
  \right)
  u_{n\mathbf k}.
  \label{eq:bloch_momentum_operator}
\end{equation}
The corresponding cell-periodic equation is
\begin{equation}
  \imath\hbar\frac{\partial}{\partial t}
  u_{n\mathbf k}(\mathbf r,t)
  =
  \hat H_{\mathbf k}^{V}(t)
  u_{n\mathbf k}(\mathbf r,t),
  \label{eq:tdks_cell_periodic}
\end{equation}
with
\begin{equation}
  \hat H_{\mathbf k}^{V}(t)
  =
  \frac{1}{2m_e}
  \left[
    -\imath\hbar\nabla
    +\hbar\mathbf k
    -
    \frac{q_e}{c}\mathbf A(t)
  \right]^2
  + \hat V_{\rm loc}[\rho](t)
  + \hat V_{{\rm NL},\mathbf k}^{V}(t).
  \label{eq:velocity_gauge_hamiltonian_k}
\end{equation}
Equivalently, using $q_e=-e$, the kinetic operator contains
$-\imath\hbar\nabla+\hbar\mathbf k+e\mathbf A(t)/c$.
Equation~\eqref{eq:velocity_gauge_hamiltonian_k} is the periodic
counterpart of the finite-system length-gauge RT-TDDFT formulation used
in our previous work,\cite{jakowski-rmg-tddft-2025} but with the position-dependent
dipole perturbation replaced by a periodic velocity-gauge coupling.

Because the dipole moment is not a well-defined bulk observable under
periodic boundary conditions, the natural time-dependent observable for
periodic optical response is the macroscopic current rather than the
molecular dipole moment.  The construction of the current operator,
including the nonlocal pseudopotential contribution, is given in
Sec.~\ref{sec:current-nlpp}; the extraction of optical response functions
from the real-time current is described in Sec.~\ref{sec:optical-response}.
